% Transcription of folder 147, batch 07 (archivists' pages 121-140).
% Pass: Opus 5.5 (claude-opus-5-5), 2026-10-09 — first pass, unchecked against the pages by a human.
%
% Fast working hand; formulas carry better than the connective prose.
% Author's pagination 60-69 on the odd sheets (121 = 60, ..., 139 = 69);
% equations numbered (185)-(234). Notation fixed for the batch:
% \mathfrak{S}_d symmetric group, \mathbb{G}_m, G_d torus, T_d local system,
% P_d torsor, D_d^! and D_d^* as written, \Pi and \widetilde{\Pi} groupoids,
% \Pi_1 fundamental groupoid, \Delta_r = [0,r] \cap \mathbb{N}.
\documentclass[11pt,a4paper]{article}
\input{../preamble/grothendieck}

\folder{147}
\batch{7}
\pages{121}{140}
\foldertitle{Structure à l'infini des Mg,ν [pages 1 à 90, dont table des matières] : notes manuscrites (s.d.).}
\dating{s.d. — le groupe « Autour de La "Longue Marche" à travers la théorie de Galois » (141 à 149) est daté [à partir de 1978]-1983}
\watermark{Édition de démonstration}

\begin{document}

\page{121}
\note{p. 60 de l'auteur. La page s'ouvre au milieu d'une construction commencée avant le lot : $D_d$, $D^{!}_d$, $T_d$, $\mathfrak{S}_d$ sont introduits plus haut.}
On désigne par $T^{!}_d$ l'image inverse de $T_d$ sur $D^{!}_d$, sur laquelle
$\mathfrak{S}_d$ opère donc (c'est un \struck{\ill{}} syst. local sur
$(D^{!}_d, \mathfrak{S}_d) \simeq D_d$). Or sur $D^{!}_d$ on a les
\add{$d$} morphismes
\[
  (185) \qquad s_i : D^{!}_d \longrightarrow D_1 \qquad (1 \leqslant i \leqslant d)
\]
permutés par $\mathfrak{S}_d$, et d'autre part, si $D_1 \to D_0$ est une
immersion \emph{régulière} \add{de codim.~1}, on a \struck{la sym\ill{} \ill{} $D_1$ \ill{}}
le faisceau \emph{normal} par l'immersion
\[
  (186) \qquad L_1 \ \text{faisceau inversible canonique sur } D_1
\]
et on peut considérer sur $D^{!}_d$ les $d$ faisceaux inv.
\[
  (187) \qquad L_{d,i} = s_i^{*}(L_1) \qquad 1 \leqslant i \leqslant d .
\]
Ils correspondent à des torseurs sous $(\mathbb{G}_m)_{D_d}$
\[
  (188) \qquad \underline{L}^{*}_{d,i} = \mathbb{G}_{m,D_d}\text{-torseur sur } D_d
  \ \text{associé au faisceau inv. } L_{d,i}
\]
et on peut considérer
\[
  (189) \qquad P^{!}_d = \prod_{1 \leqslant i \leqslant d} \underline{L}^{*}_{d,i}
  \qquad \text{torseur sous } \mathbb{G}_m^{d} .
\]
Comme $\mathfrak{S}_d$ opère sur $D^{!}_d$ permutant les \add{$s_i$ donc les}
$\underline{L}_{d,i}$ entre eux, il opère sur $P_d$, \struck{qui} \add{de}
\struck{devient} façon compatible avec son opération sur

\page{122}
\[
  (190) \qquad G^{!}_d \overset{\text{déf}}{=} \mathbb{G}_{m,D^{!}_d} \otimes_{\mathbb{Z}} T^{!}_d
  \simeq \underbrace{\mathbb{G}^{d}_{m,D^{!}_d}}
\]
\marginal{avec opération de $\mathfrak{S}_d$ \uncertain{donnée} par permutation des facteurs}
et par descente on trouve sur $D_d$ un torseur
\[
  (191) \qquad P_d \ \text{torseur sur } D_d, \text{ de groupe } G_d ,
\]
et par suite on a aussi
\[
  (192) \qquad T_d \simeq \text{syst. local des } H_1 \text{ (ou des } \pi_1)
  \text{ des fibres de } P_d .
\]

On considère le groupoïde fondamental $\Pi P_d$ de $P_d$, et
\struck{le groupoïde} \add{l'ouvert} $D^{*}_d$ est défini comme
$D_d \setminus{}$(images des $D_{d',d} \to$ \struck{$D_d$} pour $d' > d$, cf
plus bas), et plus essentiellement par la suite,
\[
  \Pi\bigl(\underbrace{P_d \mid D^{*}_d}_{\overset{\text{déf}}{=} P^{*}_d}\bigr) = \Pi P^{*}_d .
\]
\note{la lettre écrite sous l'accolade ressemble à un $T$ ; le sens demande $P^{*}_d$.}

On va considérer la variance des objets précédents $T_d$, $G_d$,
\add{$P_d$, $P^{*}_d$,} $T^{!}_d$, $G^{!}_d$, $P^{!}_d$, $\Pi P_d$, $\Pi P^{*}_d$
pour $d$ variable. Considérons donc \add{un couple}
\struck{$d \geqslant d'$}
\[
  d \geqslant d' \geqslant 0
\]

\begin{tikzcd}[column sep=small]
  & D_{d,d'} \arrow[dl, "\text{rev. ét.}"'] \arrow[dr, "\text{imm. loc.}"] & \\
  D_d & & D_{d'}
\end{tikzcd}

\note{à gauche, en marge, quelques signes biffés (« $d'=$ \ill{} »). L'indice du but de gauche, $D_d$, porte un signe peu lisible.}

\page{123}
\note{p. 61 de l'auteur.}
considérons les images inverses de $T_d$ et $T_{d'}$ sur $D_{d,d'}$, qu'on va noter
\[
  (193) \quad
  \begin{cases}
    T(d,d')_d = (T_d)_{D_{d,d'}} & \text{syst. local de rang } d \text{ sur } D_{d,d'} \\
    T(d,d')_{d'} = (T_{d'})_{D_{d,d'}} & \text{syst. local de rang } d' \text{ sur } D_{d,d'}
  \end{cases}
\]
\note{devant l'accolade, un $\mathbb{Z}_{D_{d,d'}}$ biffé.}
on a alors un homom. canonique surjectif (dit « tautologique »)
\[
  (194) \qquad T(d,d')_d \longrightarrow T(d,d')_{d'}
\]
qui, par image inverse sur $D^{!}_{d,d'} = D^{!}_d \times_{D_d} D_{d'}$
\note{lire sans doute $D_{d,d'}$ comme dernier facteur ; la virgule manque.}
(cf diagramme cart.

\begin{tikzcd}[column sep=small, row sep=small]
  (D^{!}_{d'})^{(I_d)}_{D^{!}_d} \simeq D^{!}_{d,d'} \arrow[d] \arrow[rr, dashed] \arrow[dr, "\mathfrak{S}_{d}"] & & D^{!}_{d'} \arrow[dd, "\mathfrak{S}_{d'}"] \\
  D^{!}_d \arrow[dr, "\mathfrak{S}_d"'] & D_{d,d'} \arrow[d] \arrow[dr] & \\
  & D_d & D_{d'}
\end{tikzcd}

\note{diagramme (195) ; le terme en haut à gauche est d'une lecture incertaine, et l'indice de $\mathfrak{S}$ sur la flèche $D^{!}_{d,d'} \to D_{d,d'}$ est douteux (on lit « $d^{*}$ »).}
s'identifie à l'homom. de projection
\[
  (196) \qquad \mathbb{Z}^{d} \longrightarrow \mathbb{Z}^{d'}
\]
de $\mathbb{Z}^{d}$ sur un produit partiel défini par une partie de cardinal
$d'$ de l'ens. $I_d$ en exposant : $\mathbb{Z}^{d} \simeq \mathbb{Z}^{I_d}$.

\page{124}
On déduit de (194) un morphisme correspondant (qui est épi lisse)
\[
  (197) \quad
  \begin{array}{ccc}
    G(d,d')_d & \longrightarrow & G(d,d')_{d'} \\
    \| \text{déf} & & \| \text{déf} \\
    G_d \times_{D_d} D_{d,d'} & & G_{d'} \times_{D_{d'}} D_{d,d'} \\
    \simeq T(d,d')_d \otimes_{\mathbb{Z}} \mathbb{G}_m & & \simeq T(d,d')_{d'} \otimes_{\mathbb{Z}} \mathbb{G}_m
  \end{array}
\]
et on sait que l'on a un isom. correspondant
\[
  (198) \quad
  \begin{array}{ccc}
    P(d,d')_{d'} & \simeq & \underbrace{P(d,d')_d} \wedge^{G(d,d')_d} G(d,d')_{d'} \\
    \| \text{déf} & & \| \text{déf} \\
    P_d \times_{D_d} D_{d,d'} & & P_{d'} \times_{D_{d'}} D_{d,d'}
  \end{array}
\]
\note{les deux définitions semblent interverties : on attend $P(d,d')_d = P_d \times_{D_d} D_{d,d'}$ et $P(d,d')_{d'} = P_{d'} \times_{D_{d'}} D_{d,d'}$. Transcrit tel qu'écrit.}

\struck{Au niveau des groupoïdes fondamentaux, on peut dire que
$\Pi P(d,d')_{d'}$} \add{On a des relations de transitivité pour les
morphismes précédents,} relatives \struck{\ill{}} à la donnée d'un triple
\[
  d \geqslant d' \geqslant d''
\]
\uncertain{en se plaçant} sur $D(d,d',d'')$, où on a des \struck{objets}
systèmes locaux
\[
  (199) \qquad T(d,d',d'')_d,\ T(d,d',d'')_{d'},\ T(d,d',d'')_{d''},
\]
des tores correspondants
\[
  (200) \qquad G(d,d',d'')_d,\ G(d,d',d'')_{d'},\ G(d,d',d'')_{d''}
\]

\page{125}
\note{p. 62 de l'auteur.}
et les torseurs sous ceux-ci
\[
  (201) \qquad P(d,d',d'')_d,\ P(d,d',d'')_{d'},\ P(d,d',d'')_{d''},
\]
en prenant les images inverses des quantités correspondantes sur
$D_d$, $D_{d'}$, $D_{d''}$ \note{trois symboles surchargés, puis récrits avec les indices $d$, $d'$, $d''$.} par

\begin{tikzcd}[column sep=small]
  & D_{d,d',d''} \arrow[dl] \arrow[d] \arrow[dr] & \\
  D_d & D_{d'} & D_{d''}
\end{tikzcd}

Prenant aussi les images inverses, par

\begin{tikzcd}[column sep=small]
  & D_{d,d',d''} \arrow[dl] \arrow[d] \arrow[dr] & \\
  D_{d,d'} & D_{d',d''} & D_{d,d''}
\end{tikzcd}

des morphismes du type (194), (197), (198), on trouve des transitivités

\begin{tikzcd}[column sep=small, row sep=small]
  T(d,d',d'')_d \arrow[r] \arrow[rr, bend left=20] & T(d,d',d'')_{d'} \arrow[r] & T(d,d',d'')_{d''} \\
  G(d,d',d'')_d \arrow[r] \arrow[rr, bend left=20] & G(d,d',d'')_{d'} \arrow[r] & G(d,d',d'')_{d''} \\
  P(d,d',d'')_d \arrow[r] \arrow[rr, bend left=20] & P(d,d',d'')_{d'} \arrow[r] & P(d,d',d'')_{d''}
\end{tikzcd}

\note{système (202).}

Toutes ces relations sont essentiellement tautologiques. En voici une qui ne
l'est plus, et qui fait sérieusement \uncertain{appel} : \struck{toute la force de}
l'hypothèse « croisements normaux » (dans toute sa force…)

\page{126}
\struck{On a aussi de définir sur tous} \struck{$D^{*}_{d,d'} = D_{d,d'} \mid D^{*}_d$
ouvert de $D_{d,d'}$ (complément \ill{} \ill{})}

On suppose maintenant que la situation envisagée provient d'un diviseur à
croisements normaux $D$ sur $X = D_0$, en dépliant celui-ci. Alors
\add{$\forall d \in \mathbb{N}$,} la structure \add{de} dépliement induite sur
les $D_d$ (en posant $D^{\langle d\rangle}_i = D_{d+i,d}$ ---- pour se
dispenser de développer au long \uncertain{les choses}) a aussi la même
propriété — elle provient d'un diviseur à croisements normaux \add{$\Theta_d$}
sur $D_d$, dont le complémentaire \ill{} \uncertain{est noté} $D^{*}_d$.
\struck{On note}
\note{à partir d'ici, toute la fin de la page est barrée de deux longs traits obliques ; le numéro (203) est repris à la page suivante de l'auteur (page 127 du fonds).}
\struck{(203) $D^{*}_{d,d'} = D_{d,d'} \mid D^{*}_d$}
\marginal{\struck{ouvert de $D_{d,d'}$ (complémentaire d'un diviseur à croisements normaux
$\Delta(d,d')$ sur $D_{d,d'}$, image inverse de $\Delta(d)$)}}
\note{le symbole lu $\Delta$ pourrait être un $\Theta$ mal formé.}

\struck{On \struck{pose alors} \add{peut définir des} morphismes canoniques de groupoïdes}
\[
  \text{\struck{$(203) \quad \Pi_1 P^{*}_d \xrightarrow{\ r_{d,d'}\ } \Pi_1 P^{*}_{d'}$
  \quad pour $d' < d$}}
\]
\struck{« homom. de spécialisation »}

\page{127}
\note{p. 63 de l'auteur.}
Il faut introduire ensuite les espaces analytiques de drapeaux
$D_{d_0,d_1,\dots,d_r}$ pour $d_0 \geqslant d_1 \geqslant \dots \geqslant d_r \geqslant 0$,
leur réunion est notée $D_{\Delta_r}$, où $\Delta_r$ désigne l'ens. ordonné
\[
  (203) \qquad \Delta_r = [0,r] \cap \mathbb{N}
\]
ou simplexe ordonné type de dim.~$r$. Cette notation $D_{\Delta_r}$ est
compatible avec la notation $D_I$ pour $I$ ens. fini ordonné, rajoutant une
plus petite \uncertain{élément}. Alors $D_{\Delta_r}$ dépend fonctoriellement de
façon contravariante de $\Delta_r$, par les relations croissantes, entre simplexes,
donc on a un espace analytique ½-simplicial, qu'on désigne par
$\mathcal{D}_{*}$, donc
\[
  (204) \qquad \mathcal{D}_r = \text{\struck{$\mathbf{D}$}}\, D_{\Delta_r}
  = \coprod_{\substack{(d_0,\dots,d_r) \in \mathbb{N}^{r} \\ d_0 > d_1 > \dots > d_r}}
  D_{d_0,d_1,\dots,d_r}
\]
\note{au-dessus du $\mathbf{D}$ biffé, $D_{\Delta_r}$ ajouté. L'inégalité de l'indice est ici stricte, alors qu'elle est large plus haut ; et $\mathbb{N}^{r}$ plutôt que $\mathbb{N}^{r+1}$ : transcrit tel quel.}
\struck{\ill{}} La composante $D_{d_*} = D_{d_0,d_1,\dots,d_r}$ est un revêtement
étale fini de $D_{d_0}$, s'identifie à
\[
  (205) \qquad D_{d_*} = D^{!}_{d_0} \wedge^{(\mathfrak{S}_{d_0})_{D_{d_0}}}
  \underbrace{F_{d_*}(\text{\struck{$[1,d_0]$}}\, I_{d_0})}
\]
\marginal{ens. des « drapeaux de type $d_*$ » de l'ens. fini $I_{d_0} = [1,d_0] \cap \mathbb{N}$}
\note{dans la note marginale, « drapeaux » est écrit au-dessus d'un mot biffé (\ill{}).}
On désigne par $D^{*}_{d_0,d_1,\dots,d_r}$ l'image inverse de $D^{*}_{d_0}$,
qui est donc le complémentaire d'un diviseur $\Delta_{d_0 \dots d_r}$
\note{ou $\Theta_{d_0\dots d_r}$ ; la lettre est mal formée.}
du diviseur \uncertain{$\Theta_{d_0}$} sur $D_{d_0}$. On peut bien sûr plus simpl\supplied{ement}

\page{128}
ment définir tout de suite
\[
  (206) \quad
  \begin{aligned}
    \mathcal{D}^{*}_{*} &= \coprod_{d_0 \geqslant d_1 \geqslant \dots \geqslant d_r \geqslant 0} D^{*}_{d_0,d_1,\dots,d_r} \\
    &= \text{image inverse de } D^{*}_{*} = \coprod_{d \in \mathbb{N}} D^{*}_d
    \text{ par le morphisme } o : \mathcal{D}_{*} \to D_{*}
  \end{aligned}
\]
défini par les $D_{d_0,\dots,d_r} \to D_{d_0}$, i.e.\ défini par les
\add{morphismes « premier sommet »} ($\mathcal{D}_r \to D_{*}$ transposés de
$\{0\} \to \Delta_r$)
\note{à droite, un petit schéma : un point $0$ envoyé par une flèche sur le sommet $0$ d'une chaîne verticale $0 \to 1 \to 2 \to \dots \to r$.}

Nous allons définir
\[
  (207) \quad
  \begin{cases}
    P_{d_0,d_1,\dots,d_r} = \text{image inverse de } P_{d_0} \text{ sur } D_{d_0} \\
    \qquad \text{par } D_{d_0,\dots,d_r} \to D_{d_0} \\
    P^{*}_{d_0,d_1,\dots,d_r} = P_{d_0,\dots,d_r} \mid D^{*}_{d_0,\dots,d_r}
  \end{cases}
\]
\marginal{$G_{d_0}$-torseurs}
\uncertain{c'est-à-dire} des torseurs sous les groupes analytiques (tores)
\[
  (208) \quad
  \begin{cases}
    G_{d_0,\dots,d_r} = \text{image inverse de } G_{d_0} \text{ sur } D_{d_0}
    \text{ par } D_{d_0,\dots,d_r} \to D_{d_0} \\
    G^{*}_{d_0,\dots,d_r} = G_{d_0,\dots,d_r} \mid D^{*}_{d_0,\dots,d_r} .
  \end{cases}
\]
On a donc un diagramme 2-cartésien de groupoïdes

\begin{tikzcd}[column sep=small]
  \widetilde{\Pi}_{d_0 \dots d_r} \overset{\text{déf}}{=} \Pi_1 P^{*}_{d_0,\dots,d_r} \arrow[r] \arrow[d] & \Pi_1 P^{*}_{d_0} \bigl(\overset{\text{déf}}{=} \widetilde{\Pi}_{d_0}\bigr) \arrow[d] \\
  \Pi_{d_0 \dots d_r} = \Pi_1 D^{*}_{d_0,\dots,d_r} \arrow[r] & \Pi D^{*}_{d_0} \bigl(\overset{\text{déf}}{=} \Pi_{d_0}\bigr)
\end{tikzcd}

\note{diagramme (209), marqué « 2-cart. » au centre. L'astérisque de $D^{*}_{d_0,\dots,d_r}$ en bas à gauche est écrit sur un $D$ repassé, celui de $D^{*}_{d_0}$ en bas à droite n'est pas visible.}

\page{129}
\note{p. 64 de l'auteur.}
\marginal{dès \uncertain{qu'on suppose} $\pi_2$ de $D_{d_0}$ \uncertain{nul} \ill{}…}
\ill{} les flèches verticales sont des morphismes d'extensions par des noyaux
qui sont des $\mathbb{Z}^{d_0}$ tordus — en fait par les systèmes locaux
$T_{d_0}$, et son image inverse sur $D^{*}_{d_0,\dots,d_r}$ — \uncertain{posons}
donc
\[
  (210) \quad
  \begin{cases}
    T_{d_0,\dots,d_r} = \text{image inverse de } T_{d_0} \text{ (sur } D_{d_0}\text{)}
    \text{ par } D_{d_0,\dots,d_r} \to D_{d_0} \\
    \qquad \simeq \text{syst. local des } \pi_1 \text{ (ou des } H_1)
    \text{ des fibres de } P_{d_0 \dots d_r} \\
    T^{*}_{d_0 \dots d_r} = T_{d_0 \dots d_r} \mid D^{*}_{d_0,\dots,d_r} ,
  \end{cases}
\]
\marginal{$T_{d_0}$ : syst. local \uncertain{de rang} $d_0$ ($\overset{\text{loc}}{\simeq} \mathbb{Z}^{d_0}$)}
\note{dans l'accolade, entre la première et la dernière ligne, un début de ligne biffé ($T^{*}_{d_0,\dots,d_r} =$).}
\marginal{dès \uncertain{qu'on suppose} $\pi_2(D^{*}_{d_0,\dots,d_r}) = 0$…}
dès lors que $\widetilde{\Pi}_{d_0,\dots,d_r}$ est une ext. de
$\Pi_{d_0,\dots,d_r}$ par $T^{*}_{d_0,\dots,d_r}$.

On pourra bien sûr
\[
  \text{\struck{$\Pi_{\Delta_r} = \coprod \Pi_{d_0,\dots,d_r}$}}
\]
\[
  (211) \quad
  \begin{cases}
    G_{\Delta_r} = \coprod\limits_{\substack{(d_0,\dots,d_r) \in \mathbb{N}^{r+1} \\ d_0 \geqslant d_1 \geqslant \dots \geqslant d_r}} G_{d_0,\dots,d_r}
    & \text{groupe sur } D_{\Delta_r} \\
    P_{\Delta_r} = \coprod P_{d_0,\dots,d_r} & \text{torseur sous } G_{\Delta_r} \\
    T_{\Delta_r} = \coprod T_{d_0,\dots,d_r} & \text{système local des } \pi_1 \text{ ou } H_1
    \text{ de } G_{\Delta_r} \text{ ou } P_{\Delta_r} \text{ au choix}
  \end{cases}
\]
$G^{*}_{\Delta_r}$, $P^{*}_{\Delta_r}$, $T^{*}_{\Delta_r}$ leurs restrictions à $D^{*}_{\Delta_r}$
\note{sous le $\coprod$ de $G_{\Delta_r}$, une condition biffée ($d_r \geqslant \dots \geqslant d_0$).}

\page{130}
\[
  (212) \quad
  \begin{cases}
    \Pi_{\Delta_r} = \Pi_1 D^{*}_{\Delta_r} & \text{groupoïde fond. de } D^{*}_{\Delta_r} \\
    \widetilde{\Pi}_{\Delta_r} = \Pi_1 P^{*}_{\Delta_r} & \text{id.\ de } P^{*}_{\Delta_r}
  \end{cases}
\]
donc un morph.
\[
  (213) \qquad \widetilde{\Pi}_{\Delta_r} \longrightarrow \Pi_{\Delta_r}
\]
qui fait de $\widetilde{\Pi}_{\Delta_r}$ un groupoïde extension de $\Pi_{\Delta_r}$ par
$T_{\Delta_r}$. Notons « pour mémoire »
\[
  (214) \quad
  \begin{cases}
    \Pi_{\Delta_r} = \coprod\limits_{d_0 \geqslant d_1 \geqslant \dots \geqslant d_r} \Pi_{d_0,d_1,\dots,d_r} \\
    \widetilde{\Pi}_{\Delta_r} = \coprod\limits_{d_0 \geqslant \dots \geqslant d_r} \widetilde{\Pi}_{d_0,d_1,\dots,d_r}
  \end{cases}
\]
le système local $T_{\Delta_r}$ sur $\Pi_{\Delta_r}$ \struck{provient} se réduisant aux
systèmes locaux $T_{d_0 \dots d_r}$ sur les $\Pi_{d_0,\dots,d_r}$.
\note{dans la marge gauche, en face de (214), un petit trait barré sans texte lisible.}

Les $\Pi_{\Delta_r}$ n'ont pas de « dépendance simpliciale » par rapport à
$\Delta_r$, p.~ex. on n'a pas de morphisme naturel entre
$\Pi_{d,d'} = \Pi_1 D^{*}_{d,d'}$, et $\Pi_{d'} = \Pi_1 D^{*}_{d'}$, \uncertain{l'application}

\page{131}
\note{p. 65 de l'auteur.}
contre il y en a entre $\Pi_{d,d'}$ et $\Pi_d$. Précisons ce point, en considérant
un morphisme
\[
  (215) \qquad \Delta_r \xrightarrow{\ \alpha\ } \Delta_{r'} \qquad \text{tel que } \alpha(0) = 0
\]
\note{« tel que $\alpha(0)=0$ » doublement souligné.}
alors $D_{\Delta_{r'}} \to D_{\Delta_r}$ est un morphisme étale fini, donc aussi
$D^{*}_{\Delta_{r'}} \to D^{*}_{\Delta_r}$ (induit au-dessus de $D^{*}_{\Delta_r}$),
donc on trouve un diagramme 2-cartésien (qui \struck{\ill{}} \add{\ill{}}
(\struck{209}) \ill{} se déduit de (209))

\begin{tikzcd}
  \widetilde{\Pi}_{\Delta_{r'}} \arrow[r] \arrow[d] & \widetilde{\Pi}_{\Delta_r} \arrow[d] \\
  \Pi_{\Delta_{r'}} \arrow[r] & \Pi_{\Delta_r}
\end{tikzcd}

\note{diagramme (216), marqué « 2-cart. » au centre.}
\marginal{Remarquons que les $\widetilde{\Pi}_{\Delta_r} \to \Pi_{\Delta_r}$,
$\Pi_{\Delta_{r'}} \to \Pi_{\Delta_r}$ \ill{} et $\widetilde{\Pi} \to$ \ill{}
\uncertain{correspondant} $\widetilde{\Pi}_d \to \Pi_d$, \uncertain{conservent} les
\struck{\ill{}} \ill{} \ill{} $D^{*!}_d$ sur $D_d$ \ill{}}
\note{note marginale écrite en oblique et en partie raturée, dont seule une partie se lit.}
où les flèches horizontales induisent \add{\ill{}} sur les $\pi_1$ des
\uncertain{injections} : \ill{} image d'indice fini (au fond, dire que
$\Pi_{\Delta_{r'}} \to \Pi_{\Delta_r}$ fait de $\Pi_{\Delta_{r'}}$ un rev. fini de
$\Pi_{\Delta_r}$).

Toutes ces relations sont encore de nature essentiellement formelle et
triviale, et n'utilisent dans une certaine mesure \uncertain{que} ; \ill{}
l'hypothèse spécifique des croisements normaux, et \uncertain{aussi} que
\uncertain{la} \uncertain{résolution} des

\page{132}
$D_{\Delta_r}$ par les $D^{*}_{\Delta_r}$, pour définir les $\Pi_{\Delta_r}$ et
$\widetilde{\Pi}_{\Delta_r}$, \add{seules,} \struck{jouent} un rôle particulier
\add{ici}.

Pour $d_0$ \emph{fixé}, les systèmes
\[
  \widetilde{\Pi}_{d_0 d_1 \dots d_r} \longrightarrow \Pi_{d_0 \dots d_r}
\]
forment un \uncertain{tour} de groupoïdes semi-simpliciaux, mais de nature
essentiellement triviale (il faudrait cependant examiner l'occasion de
\uncertain{plus près} cette mi-distante trivialité — \uncertain{elle s'explicite}
sans peine en termes de la théorie de Galois des revêtements de $D^{*}_{d_0}$
\uncertain{resp.} de $P^{*}_{d_0}$…).

Nous allons maintenant introduire une structure de nature tout à fait
non-formelle, qui va provenir de l'hypothèse des croisements normaux
\uncertain{dans son entier}, et s'exprimer par le fait que les $\widetilde{\Pi}_{\Delta_r}$
(NB et non les $\Pi_{\Delta_r}$ !) forment un groupoïde ½-simplicial, pour
$\Delta_r$ variable. En d'autres

\page{133}
\note{p. 66 de l'auteur.}
termes, il faut introduire, \struck{\ill{}} pour un morphisme
\[
  (218) \qquad \Delta_r \xrightarrow{\ \alpha\ } \Delta_{r'}
\]
un morphisme correspondant de groupoïdes
\[
  (219) \qquad \boxed{\ \widetilde{\Pi}_{\Delta_{r'}} \xrightarrow{\ \widetilde{\alpha}^{*}\ } \widetilde{\Pi}_{\Delta_r}\ }
\]
\note{la formule est encadrée ; au but on lit $\Pi_{\Delta_r}$ sans tilde visible.}
qu'on appellera l'\emph{homom.\ de spécialisation} associé à $\alpha$.
\marginal{double trait vertical en marge de ces lignes}
Lorsque $\alpha$ se factorise canoniquement en un homom.\ surjectif suivi d'un
homom.\ injectif de simplexes, \ill{} \uncertain{il est} nécessaire d'examiner
séparément les deux cas. Mais si $\alpha$ est surjectif
\add{on a $\alpha(0) = 0$} (« homom.\ de dégénérescence »),
\struck{l'homom.\ morphisme} \add{dans ce cas dans le cas déjà envisagé —}
\struck{correspondant} \add{\uncertain{donc} \ill{}} dans ce cas
$D^{*}_{\Delta_{r'}} \to D^{*}_{\Delta_r}$ \add{est} \ill{} \uncertain{décomposé}
en des morphismes
\[
  D^{*}_{d'_0,\dots,d'_{r'}} \longrightarrow D^{*}_{d_0,\dots,d_r}
\]
où $(d_0,\dots,d_r)$ est déduit de $(d'_0,\dots,d'_{r'})$ en répétant certains
de ses éléments — et la flèche (220) \uncertain{en est} un iso — dans
\emph{si $\alpha$ est surjectif}
\marginal{On voit que les \uncertain{considérants} \uncertain{des}
$\widetilde{\Pi}_{d_0 \dots d_r}$, $\Pi_{d_0,\dots,d_r}$ \ill{} \uncertain{pour}
$d_0,\dots,d_r$ distincts, \uncertain{par} (220), et \ill{} \uncertain{le} \ill{}
\emph{stricts} de \ill{} \emph{stricts} \ill{} \ill{} \uncertain{suffisent}
\ill{} \ill{}}
\note{note marginale écrite en oblique dans la marge gauche, à peine lisible ; la numérotation saute de (216) à (218), et (220) est cité avant d'être écrit.}

\page{134}
\emph{de dégénérescence}, \emph{alors} (219) est \emph{une \struck{\ill{}}
\add{bijection}, \uncertain{\ill{}}}, de nature essentiellement triviale.

On est donc réduit à regarder le cas d'une application \emph{injective} $\alpha$
(« opération face »), et le seul cas qui pose problème est celui où
$\alpha(0) \neq 0$. Donc si
\[
  d_{*} = \{d_0 > d_1 > \dots > d_r\}, \qquad \text{et si}
\]
\[
  (221) \qquad d'_{*} = \{d'_0 > d'_1 > \dots > d'_{r'}\}
\]
est \struck{\ill{}} \add{contenue dans} $d_{*}$,
$d'_0 \neq d_0$ i.e.\ $d'_0 < d_0$, il faut définir l'homom.
\[
  (222) \qquad \Pi_1 P^{*}_{d_0,d_1,\dots,d_r} \longrightarrow \Pi_1 P^{*}_{d'_0,\dots,d'_r}
\]
\marginal{attention, on est dans une situation de \ill{} \uncertain{spécialisation}
\uncertain{et} on retrouve des \ill{} \uncertain{distincts} de $d_{*}$ et $d'_{*}$ !}
\marginal{On se réduit au cas $d_0,d_1,\dots,d_r$ strictement \uncertain{décroissant}
(221)… En fait \ill{} \uncertain{il suffit de} \ill{} \ill{}
« face » \ill{} \ill{} correspondant \ill{} $\alpha(0) \neq 0$ (222)
$D_{d_0,d_1,\dots,d_r} \to D_{d_1,\dots,d_r}$}
\note{deux notes marginales, l'une à droite, entourée, l'autre en oblique dans la marge gauche ; lecture très partielle.}
Notons que l'on a
\[
  (223) \qquad D^{*}_{d_0,\dots,d_r} \longrightarrow D_{d'_0,\dots,d'_{r'}}
\]
(mais \emph{pas} un morphisme entre les espaces avec $*$ !), qui est une
immersion \add{locale} d'espaces analytiques lisses, de codim.\ $d_0 - d'_0$,
qui \add{est \uncertain{incluse}} \ill{} dans le diviseur

\page{135}
\note{p. 67 de l'auteur.}

\begin{tikzcd}
  D_{d_0,\dots,d_r} \arrow[r] \arrow[d] & D_{d'_0 \dots d'_{r'}} \arrow[d] \\
  D_{d_0,d'_0} \arrow[r] & D^{!}_{d'_0}
\end{tikzcd}

\note{diagramme (224) ; le numéro est écrit sur un autre chiffre. La flèche verticale de droite est un simple trait sans pointe visible ; l'exposant du coin inférieur droit est surchargé ($*$ puis $!$).}
où les flèches verticales sont étales, et la 2\textsuperscript{e} flèche horizontale est
une immersion locale de codim.\ $d_0$ \struck{$d'_0$} \struck{\ill{}}
\[
  (225) \qquad d = d_0 - d'_0
\]
\note{la lettre de gauche est surchargée ; le numéro (225) est repris plus bas pour une autre formule.}
[\struck{D}ans le typique (auquel on se réduit) où $d'_{*} = \{d_1 > \dots > d_r\}$,
le diagramme (224) est cartésien — \add{plus généralement} \struck{\ill{}}, il
\struck{\ill{}} \struck{\ill{}} cartésien, \add{si $d'_0 = d_1$} et c'est trivial.]
\marginal{\emph{NB} Dans le cas général, $D_{d_*} \xrightarrow{\text{fini étale}} D^{*!}_{d_*} \to D_{d_0 d'_0} \to D^{!}_{d'_0}$ et $D^{*!}_{d_*} \to D^{!*}_{d'_*} \xrightarrow{\text{cart}} D^{!}_{d'_0}$, avec des flèches pointillées $D_{d_*} \to D^{!}_{d'_0}$ et $D_{d_0 d'_0} \to D^{!}_{d'_0}$}
\note{petit diagramme écrit verticalement dans la marge gauche ; les exposants $*$ et $!$ y sont d'une lecture incertaine, et la description qu'on en donne ici est approximative.}
Considérons le morphisme induit
\[
  (225) \qquad D^{*}_{d_{*}} \longrightarrow D_{d'_{*}}
\]
(attention, \uncertain{mais} \uncertain{non} $*$ au \uncertain{départ}
\uncertain{et} $D_{d_*}$ seulement, mais : $D_{d'_*}$ !).
Nous munissons $D_{d'_*}$ du diviseur $\Theta_{d'_*}$ image inverse de $\Theta_{d_0}$
sur $D_{d_0}$. On désigne par $\Theta^{(p)}_{d'_*}$ le « squelette de dim.~$p$ » de
$\Theta^{p}_{d'_*}$, et \struck{\ill{}} on pose

\page{136}
\[
  (226) \qquad \Theta^{d*}_{d'_*} \overset{\text{déf}}{=} \Theta^{d}_{d'_*} \setminus \Theta^{d+1}_{d'_*}
\]
\note{les exposants $d$ sont écrits sur d'autres signes ($d_*$ ?).}
squelette « privé » de codim.~$d$ du $\Theta_{d'_*}$ dans $D_{d'_*}$.

Ceci posé, on sait que 225 induit un \struck{morphisme} isomorphisme \emph{local}
\emph{de} $D_{d_*}$ \emph{sur} \struck{\ill{}} $\Theta^{d*}_{d'_*}$. Plus
précisément, \struck{$\Theta^{d*}_{d'_*}$ s'identifie à :}
\struck{$D_{(d_0,d'_0,\dots,d'_r)} : D_{\tilde{d}'_*}$, où}
\[
  (227) \qquad \Theta^{(d)}_{d'_*} \xleftarrow{\ \sim\ } D_{\tilde{d}'_*}
  \qquad \text{où } \tilde{d}'_* = (d_0, d'_0, \dots, d'_r)
\]
\note{le numéro est écrit sur un autre chiffre.}
et cet isom.\ induit un isom.
\[
  (228) \qquad D^{*}_{\tilde{d}'_*} \xrightarrow{\ \sim\ } \Theta^{(d)*}_{d'_*} ,
\]
d'autre part (223) $D_{d_*} \to D_{d'_*}$ se factorise en
\[
  (229) \qquad D_{d_*} \longrightarrow D_{\tilde{d}'_*} \simeq \Theta^{(d)*}_{d'_*} \longrightarrow D_{d'_*}
\]
\note{l'isomorphisme $D_{\tilde{d}'_*} \simeq \Theta^{(d)*}_{d'_*}$ est écrit verticalement sous le terme du milieu ; l'astérisque de $\Theta^{(d)*}$ y est repassé.}
donc c'est un morphisme étale fini, suivi \add{\uncertain{ensuite}} d'un morphisme
d'inclusion de $\Theta^{(d)}_{d'_*}$ dans $D_{d'_*}$. De plus, comme $D^{*}_{d_*}$ est

\page{137}
\note{p. 68 de l'auteur.}
évidemment l'image inverse de $D^{*}_{\tilde{d}'_*}$, i.e.\ \uncertain{aussi} l'image
inverse de $\Theta^{(d)*}_{d'_*}$.

\struck{Considérons un diagramme} \uncertain{Simplifions}, \ill{} on est ramené
à la \emph{situation suivante} : \struck{$X$} \add{$Y$} espace analytique lisse,
muni d'un diviseur $\Theta$ à croisements normaux, $D''$ et $D'$ deux espaces
analytiques (lisses), et
\[
  (230) \qquad D''^{*} \longrightarrow D' \longrightarrow \text{\struck{$X$}}\ Y
\]
\note{l'astérisque de $D''^{*}$ est écrit sur un autre signe, et le $Y$ final remplace un $X$ biffé.}
des morphismes, qui sont \emph{localement} des immersions, on suppose $D''$
de codim.\ $d_0$, $D'$ de codim.\ $d_1$ dans $Y$, donc $D''$ de codim.\
$d = d_1 - d_0$ dans $D'$, et que $D' \to Y$ se factorise par $\Theta^{(d_1)}$
(squelette de dim.~$d_1$ de $\Theta$),
\struck{\ill{}} \add{$d_1 > d_0 \geqslant 0$}
et $D''^{*} \to Y$ se factorise par $\Theta^{(d_0)*}$ (squelette
\struck{\ill{}} de codim.\ $d_0$ de $\Theta$). On désigne par \struck{\ill{}}
$D''^{*}$ et $D'^{*}$ les images \add{inverses} \struck{\ill{}} de $\Theta^{(d_0)*}$
et $\Theta^{(d_1)*}$ \struck{\ill{}}. \add{\ill{}} donc $\Theta$ \ill{}
une obstruction \ill{} constituée le long de $D''^{*}$, \ill{} \uncertain{réunion} locale
de $d_0$ branches non singulières \uncertain{de codim.} $d_0$ qui se coupent
suivant $D''$). On considère les « fibrés principaux multinormaux »
$P_{D''}$, $P_{D'}$, associés aux immersions locales $D'' \to Y$, $D' \to Y$,
et les groupoïdes
\marginal{$d_1 > d_0 \geqslant 0$ (cas $d_0 = 0$ \uncertain{non} exclus, \ill{} \ill{} \ill{} trivial) \ill{} \ill{} \uncertain{cartésien} ; petit diagramme : $D''^{*} \to D'^{*} \dashrightarrow D_0 = Y$, $D''^{*} \xrightarrow{\text{ét.}} D_{d_0,d_1} \to D_{d_1} \to D_0$, marqué « cart. »}
\note{note marginale écrite en oblique dans la marge gauche ; la description du diagramme est approximative.}

\page{138}
fondamentaux $\Pi_1 P_{D''^{*}}$, $\Pi_1 P_{D'^{*}}$,
\marginal{\struck{où $D'^{*}$ est l'image inverse de $\Theta^{(d_1)*}$}}
\note{note biffée et encadrée au-dessus de la ligne.}
Ceci posé, on a un morphisme canonique (dit « de spécialisation »)
\[
  (231) \qquad \boxed{\ \Pi_1 P_{D''^{*}} \xrightarrow{\ r_{D'',D'}\ } \Pi_1 P_{D'^{*}}\ }
\]
\struck{compatible avec les opérations de} \add{\struck{\uncertain{principalement}}}
On donnera la construction de ce morphisme, et ses principales propriétés
formelles, dans la suite. Notons dès à présent que \uncertain{\ill{}} un
morphisme trivial, déduit de l'inclusion
\struck{$P_{D'^{*}} = P_{D'} \mid D'^{*} \to P_{D'}$}
\[
  P_{D'^{*}} = P_{D'} \mid D'^{*} \hookrightarrow P_{D'}
\]
\[
  (232) \qquad \Pi_1 P_{D'^{*}} \longrightarrow \Pi_1 P_{D'}
\]
donc le composé de (231) et (232) fournit
\[
  (233) \qquad \Pi_1 P_{D'^{*}} \longrightarrow \Pi_1 P_{D'} ,
\]
\note{on attend $\Pi_1 P_{D''^{*}}$ comme source du composé ; la page porte $D'^{*}$, transcrit tel quel.}
qui, lui, est de nature essentiellement triviale, car on a en effet un
morphisme évident
\[
  P_{D'} \longrightarrow P_{D''} .
\]
En fait,

\page{139}
\note{p. 69 de l'auteur.}
on doit avoir commutativité dans

\begin{tikzcd}
  \Pi_1 P_{D''^{*}} \arrow[r, "r_{D'',D'}"] \arrow[d, "i_{D''}"'] & \Pi_1 P_{D'^{*}} \arrow[d, "i_{D'}"] \\
  \Pi_1 P_{D''} \arrow[r] & \Pi_1 P_{D'}
\end{tikzcd}

\note{diagramme (234).}
On fera la distinction entre la différence entre la flèche composée
$i_{D'} r_{D'',D'}$, de nature triviale, et la flèche $r_{D'',D'}$, qui ne
correspond à aucun morphisme $P_{D''^{*}} \to P_{D'^{*}}$ (\emph{NB} l'image de
$D''$ dans $D'$ est contenue justement dans $D' \setminus D'^{*}$ !).

Le principe de la construction de (231) est le suivant. Pour une immersion
locale d'espaces analytiques lisses telle que $D' \to Y$ (\uncertain{condition} :
un diviseur $\Theta$ à croisements normaux sur $Y$, i.e.\ se factorisant par
$D' \xrightarrow{\text{ét.}} D \to Y = D_0$),\note{avant $D$, un mot biffé illisible ; l'indice de $D$ est lui aussi biffé et illisible.}
on introduit \struck{\ill{}} (en fait, indépendamment de la donnée de $\Theta$)
la « voisinage tubulaire de $D'$ \struck{\ill{}} dans $Y$ », comme un germe
d'espace analytique \struck{\ill{}} \uncertain{autour} de $D'$

\page{140}
p.~ex.\ si l'image de $D'$ dans $Y$ est figurée par
\drawing{une courbe formant une boucle qui se recoupe en un point double marqué d'un gros point.}
le voisinage tubulaire \add{dans $Y$} en question sera figuré par
\drawing{la même boucle épaissie en une bande hachurée transversalement, les deux branches de la bande se croisant au point double.}
et \uncertain{non} dans $Y$ par
\drawing{un arc de bande hachurée en forme de U, non recoupé, avec une flèche marquée $D'$ désignant l'âme de la bande.}
\note{la parenthèse fermante placée au-dessus de la troisième figure ne s'ouvre nulle part sur la page.}

Cette construction est possible dans le cadre (p.~ex.) des
\add{\uncertain{prémulti}. (\uncertain{locales}) de} variétés différentiables,
\add{\uncertain{paracompactes}}, en tant que germe de variété autour de $D'$
(défini : \uncertain{à} \uncertain{unique} \uncertain{près}…). Dans le cadre des
\emph{multiplicités} \struck{topologiques} \add{diff.\ (voire topologiques)}, on
peut encore définir \add{un voisinage} \struck{\ill{}} \add{tubulaire} comme un
germe de \struck{topos,} \add{multiplicité diff.}\ \uncertain{abstraite} :
\uncertain{partie} d'une situation d'\struck{espaces analytiques} \add{variétés
diff.\ (\uncertain{ou top.})} paracompactes. On peut aussi, \uncertain{il}
\uncertain{semble}-t-il, \uncertain{donner}
\note{la page, et le lot, s'arrêtent au milieu de la phrase ; l'argument se poursuit au-delà du lot.}

\end{document}
